

\medskip
\noindent\SGAInlineRefTarget{sga3:iii:paragraph:4-11}\textbf{4.11.}
We now propose to study the following situation. Let \(S\),
\(S_{\mathscr J}\), and \(S_0\) be as in \SGARefLink{sga3:iii:proposition:4-8}{4.8}. We have three
\(S\)-schemes \(X,X',T\), a subscheme \(Y\) of \(X\) (respectively a
subscheme \(Y'\) of \(X'\)), and morphisms \(f:T\to X'\) and
\(g:X'\to X\).

\begin{center}
  \SGAThreeNativeInput[0.62\textwidth]
    {native_diagrams/exp3/diagram_043_exp3_p055_411_setup.tex}
\end{center}

We suppose that, after reduction modulo \(\mathscr J\), this diagram can
be completed to a commutative diagram

\begin{center}
  \SGAThreeNativeInput[0.72\textwidth]
    {native_diagrams/exp3/diagram_044_exp3_p056_411_reduction.tex}
\end{center}

By \SGARefLink{sga3:iii:proposition:4-8}{4.8} we therefore have the obstructions
\begin{align*}
 c(X,Y,g\circ i')&\in
 \operatorname{Hom}_{\mathscr O_{Y'_0}}
 \bigl(i_0'^*g_0^*(\mathscr N_{Y/X}
   \otimes_{\mathscr O_Y}\mathscr O_{Y_0}),
   \mathscr J\mathscr O_{Y'}\bigr),\\
 c(X',Y',f)&\in
 \operatorname{Hom}_{\mathscr O_{T_0}}
 \bigl(f_0^*(\mathscr N_{Y'/X'}
   \otimes_{\mathscr O_{Y'}}\mathscr O_{Y'_0}),
   \mathscr J\mathscr O_T\bigr),\\
 c(X,Y,g\circ f)&\in
 \operatorname{Hom}_{\mathscr O_{T_0}}
 \bigl(f_0^*g_0^*(\mathscr N_{Y/X}
   \otimes_{\mathscr O_Y}\mathscr O_{Y_0}),
   \mathscr J\mathscr O_T\bigr),
\end{align*}
whose relations we seek to compute.

\medskip
\noindent\SGAInlineRefTarget{sga3:iii:lemma:4-12}\textbf{Lemma 4.12.}
\textit{Suppose that \(Y'\) is flat over \(S\), so that
\[
 \mathscr J\otimes_{\mathscr O_{S_0}}\mathscr O_{Y'_0}
 =\mathscr J\mathscr O_{Y'}.
\]
Then:}
\begin{enumerate}
\item[(i)] \textit{There is a natural morphism}
\begin{align*}
 b_{f_0}:&\
 \operatorname{Hom}_{\mathscr O_{Y'_0}}
 \bigl(i_0'^*g_0^*(\mathscr N_{Y/X}
   \otimes_{\mathscr O_Y}\mathscr O_{Y_0}),
   \mathscr J\mathscr O_{Y'}\bigr)\\
 &\longrightarrow
 \operatorname{Hom}_{\mathscr O_{T_0}}
 \bigl(f_0^*g_0^*(\mathscr N_{Y/X}
   \otimes_{\mathscr O_Y}\mathscr O_{Y_0}),
   \mathscr J\mathscr O_T\bigr).
\end{align*}
\item[(ii)] \textit{There is also a natural morphism, functorial in
\(T\),}
\footnote{Editor's note: We have removed the hypothesis that \(Y'_0\)
be locally complete intersection in \(X'_0\), which is superfluous by
\SGARefLink{sga3:iii:remark:4-8-0}{4.8.0}. We have also added that \(a_{g_0}(f_0)\) is
\lq\lq functorial in \(T\)\rq\rq, a fact that plays a crucial role in
the proof of \SGARefLink{sga3:iii:lemma:4-17}{4.17}.}
\begin{align*}
 a_{g_0}(f_0):&\
 \operatorname{Hom}_{\mathscr O_{T_0}}
 \bigl(f_0^*(\mathscr N_{Y'/X'}
   \otimes_{\mathscr O_{Y'}}\mathscr O_{Y'_0}),
   \mathscr J\mathscr O_T\bigr)\\
 &\longrightarrow
 \operatorname{Hom}_{\mathscr O_{T_0}}
 \bigl(f_0^*g_0^*(\mathscr N_{Y/X}
   \otimes_{\mathscr O_Y}\mathscr O_{Y_0}),
   \mathscr J\mathscr O_T\bigr).
\end{align*}
\end{enumerate}

\noindent\textit{Proof.}
\footnote{Editor's note: The proof has been expanded to make the
\lq\lq functoriality in \(T\)\rq\rq\ of \(a_{g_0}\) explicit.}
First note that, with \(X,X',Y,Y'\) fixed, giving a \(T\) as above is
equivalent to giving a morphism
\((f,f_{\mathscr J}):T\to X'\times_{X'^{+}}Y'^{+}\). Put
\(Z=X'\times_{X'^{+}}Y'^{+}\), and let \(M\) and \(M'\) be the
\(Z\)-functors defined, for every \(f:T\to Z\), by
\begin{align*}
 M(T)&=
 \operatorname{Hom}_{\mathscr O_{T_0}}
 \bigl(f_0^*g_0^*(\mathscr N_{Y/X}
   \otimes_{\mathscr O_Y}\mathscr O_{Y_0}),
   \mathscr J\mathscr O_T\bigr),\\
 M'(T)&=
 \operatorname{Hom}_{\mathscr O_{T_0}}
 \bigl(f_0^*(\mathscr N_{Y'/X'}
   \otimes_{\mathscr O_{Y'}}\mathscr O_{Y'_0}),
   \mathscr J\mathscr O_T\bigr).
\end{align*}

In all cases there is a commutative diagram

\begin{center}
  \SGAThreeNativeInput[0.65\textwidth]
    {native_diagrams/exp3/diagram_045_exp3_p057_412_tensor_square.tex}
\end{center}

Since \(Y'\) is flat over \(S\), the left vertical arrow is an
isomorphism, and hence we obtain a morphism of
\(\mathscr O_{T_0}\)-modules
\[
 f_0^*(\mathscr J\mathscr O_{Y'})\longrightarrow
 \mathscr J\mathscr O_T.
\]
It induces a morphism of abelian groups
\begin{multline*}
 \operatorname{Hom}_{\mathscr O_{T_0}}
 \bigl(f_0^*g_0^*(\mathscr N_{Y/X}
  \otimes_{\mathscr O_Y}\mathscr O_{Y_0}),
  f_0^*(\mathscr J\mathscr O_{Y'})\bigr)\\
 \longrightarrow M(T),
\end{multline*}
and, on composing with the morphism
\begin{multline*}
 M(Y')\longrightarrow
 \operatorname{Hom}_{\mathscr O_{T_0}}
 \bigl(f_0^*g_0^*(\mathscr N_{Y/X}
  \otimes_{\mathscr O_Y}\mathscr O_{Y_0}),
  f_0^*(\mathscr J\mathscr O_{Y'})\bigr)
\end{multline*}
induced by \(f_0^*\), we obtain the morphism
\(b_{f_0}:M(Y')\to M(T)\).

Likewise, in all cases there is a diagram

\begin{center}
  \SGAThreeNativeInput[0.67\textwidth]
    {native_diagrams/exp3/diagram_046_exp3_p057_412_conormal_square.tex}
\end{center}

Since \(Y'\) is flat over \(S\), the second vertical arrow is an
isomorphism by \SGARefLink{sga3:iii:remark:4-8-0:equation:4-8-0}{4.8.0}. We therefore obtain an
\(\mathscr O_{Y'_0}\)-morphism
\[
 i_0'^*g_0^*(\mathscr N_{Y/X}
   \otimes_{\mathscr O_Y}\mathscr O_{Y_0})
 \longrightarrow
 \mathscr N_{Y'/X'}
   \otimes_{\mathscr O_{Y'}}\mathscr O_{Y'_0}.
\]
This induces a morphism
\(a_{g_0}(\operatorname{id}_{Y'}):M'(Y')\to M(Y')\) and, for every
\(f:T\to Z\), a morphism \(a_{g_0}(f):M'(T)\to M(T)\), such that the
following diagram commutes:

\begin{center}
  \SGAThreeNativeInput[0.50\textwidth]
    {native_diagrams/exp3/diagram_047_exp3_p057_412_functor_square.tex}
\end{center}

Here \(b'_{f_0}\) is defined in the same way as \(b_{f_0}\).
\hfill\textsc{q.e.d.}

\medskip
\noindent\SGAInlineRefTarget{sga3:iii:remark:4-12-1}\textbf{Remark 4.12.1.}
\footnote{Editor's note: This remark has been added for use in the
proof of \SGARefLink{sga3:iii:lemma:4-17}{4.17}.}
\textit{Let \(M_0\) and \(M'_0\) be the \(Y'_0\)-functors defined, for
every \(f:T_0\to Y'_0\), by}
\begin{align*}
 M_0(T)&=
 \operatorname{Hom}_{\mathscr O_{T_0}}
 \bigl(f_0^*g_0^*(\mathscr N_{Y/X}
  \otimes_{\mathscr O_Y}\mathscr O_{Y_0}),
  \mathscr J\otimes_{\mathscr O_{S_0}}\mathscr O_{T_0}\bigr),\\
 M'_0(T)&=
 \operatorname{Hom}_{\mathscr O_{T_0}}
 \bigl(f_0^*(\mathscr N_{Y'/X'}
  \otimes_{\mathscr O_{Y'}}\mathscr O_{Y'_0}),
 \mathscr J\otimes_{\mathscr O_{S_0}}\mathscr O_{T_0}\bigr).
\end{align*}

Note at once that \(Z_0=Y'_0\), and that on the category of
\(Z\)-schemes \(T\) flat over \(S\), the functors \(M\) and \(M'\)
coincide respectively with
\(\prod_{S_0/S}M_0\) and \(\prod_{S_0/S}M'_0\). In that case,
\(b_{f_0}\) is simply the morphism
\[
 f_0^*:M_0(Y'_0)\longrightarrow M(T_0)
\]
induced by \(f_0\). Moreover, for every morphism \(u:U\to T\), if
\(h=f\circ u\), the following diagram commutes:

\begin{center}
  \SGAThreeNativeInput[0.50\textwidth]
    {native_diagrams/exp3/diagram_048_exp3_p058_4121_naturality_square.tex}
\end{center}

Thus \(a_{g_0}\) becomes a morphism of functors
\[
 \prod_{S_0/S}M'_0\longrightarrow\prod_{S_0/S}M_0.
\]

\medskip
\noindent\SGAInlineRefTarget{sga3:iii:proposition:4-13}\textbf{Proposition 4.13.}
\textit{Suppose that \(Y'\) is flat over \(S\). Then}
\[
 c(X,Y,g\circ f)=a_{g_0}\bigl(c(X',Y',f)\bigr)
   +b_{f_0}\bigl(c(X,Y,g\circ i')\bigr).
\]

Since the definitions of the various obstructions and of the morphisms
\(a_{g_0}\) and \(b_{f_0}\) are local, it plainly suffices to verify the
formula when all the schemes in question are affine. Write
\[
 \begin{gathered}
 S=\operatorname{Spec}(\Lambda),\qquad
 S_{\mathscr J}=\operatorname{Spec}(\Lambda/\mathscr J),\qquad
 S_0=\operatorname{Spec}(\Lambda/\mathscr I),\\
 T=\operatorname{Spec}(C),\qquad X=\operatorname{Spec}(A),\qquad
 Y=\operatorname{Spec}(A/I_Y)=\operatorname{Spec}(B),\\
 X'=\operatorname{Spec}(A'),\qquad
 Y'=\operatorname{Spec}(A'/I_{Y'})=\operatorname{Spec}(B').
 \end{gathered}
\]
There is therefore a diagram of rings and ideals

\begin{center}
  \SGAThreeNativeInput[0.48\textwidth]
    {native_diagrams/exp3/diagram_049_exp3_p058_413_ring_ideal_diagram.tex}
\end{center}

Let us study the terms of the formula to be proved. In what follows,
for \(x\in I_Y\) (respectively \(u\in I_{Y'}\)), write \(\bar x\)
(respectively \(\bar u\)) for its image in \(I_Y/I_Y^2\)
(respectively \(I_{Y'}/I_{Y'}^2\)). If \(m\) belongs to a
\(\Lambda\)-module \(M\), write \(\bar m_0\) for its image in
\(M_0=M/\mathscr I M\).

We saw that \(c=c(X,Y,g\circ f)\) is the unique \(C_0\)-morphism
\[
 I_Y/I_Y^2\otimes_B C_0\longrightarrow\mathscr J C
\]
such that \(c(\bar x\otimes1)=f(g(x))\) for every \(x\in I_Y\).

Fix \(x\in I_Y\). Since \(g_{\mathscr J}(Y'_{\mathscr J})\subset
Y_{\mathscr J}\), one has \(g(x)\in I_{Y'}+\mathscr J A'\). Write
\[
 g(x)=x'+\sum_i\lambda_i a'_i,
 \qquad x'\in I_{Y'},\quad \lambda_i\in\mathscr J,\quad a'_i\in A'.
\]
Then
\begin{equation*}
 \tag{1}\label{sga3:iii:proposition:4-13:equation:1}
 c(X,Y,g\circ f)(\bar x\otimes1)
 =f(g(x))=f(x')+\sum_i\lambda_i f(a'_i).
\end{equation*}

Consider \(a_{g_0}(c(X',Y',f))\). By the definitions just given, it is
defined by the following diagram:

\begin{center}
  \SGAThreeNativeInput[0.86\textwidth]
    {native_diagrams/exp3/diagram_050_exp3_p059_413_ag0_obstruction_diagram.tex}
\end{center}

Thus
\(a_{g_0}(c(X',Y',f))(\bar x\otimes1)=f(u)\), where \(u\in I_{Y'}\)
has image \(\bar u\) in \(I_{Y'_0}/I_{Y'_0}^2\) satisfying
\[
 \bar u_0\otimes1=\bar g_0(\bar x_0)\otimes1
 =g_0(\bar x_0)\otimes1.
\]
We may therefore take \(u=x'\), which gives
\begin{equation*}
 \tag{2}\label{sga3:iii:proposition:4-13:equation:2}
 a_{g_0}(c(X',Y',f))(\bar x\otimes1)=f(x').
\end{equation*}

Finally consider \(b_{f_0}(c(X,Y,g\circ i'))\). By hypothesis, the
morphism of \(\Lambda_0\)-algebras \(f_0:A'_0\to C_0\) factors through
\(B'_0\). Since \(Y'\) is flat, and hence
\[
 \mathscr J\otimes_{\Lambda_0}B'_0\xrightarrow{\sim}\mathscr J B',
\]
we obtain a morphism of \(B'_0\)-modules
\(\psi:\mathscr J B'\to\mathscr J C\) making the following diagram
commute:

\begin{center}
  \SGAThreeNativeInput[0.74\textwidth]
    {native_diagrams/exp3/diagram_051_exp3_p059_413_tensor_J_square.tex}
\end{center}

Let
\[
 \varphi:\mathscr J B'\otimes_{B'_0}C_0\longrightarrow\mathscr J C
\]
be the morphism of \(C_0\)-modules induced by \(\psi\). Then, for every
\(a'\in A'\) and \(\lambda\in\mathscr J\),
\begin{equation*}
\tag{\(\dagger\)}\label{sga3:iii:proposition:4-13:equation:dagger}
 \varphi(\lambda\pi'(a')\otimes1)=\lambda f(a').
\end{equation*}

The element \(b_{f_0}(c(X,Y,g\circ i'))\) is then defined by the
commutative diagram

\begin{center}
  \SGAThreeNativeInput[0.62\textwidth]
    {native_diagrams/exp3/diagram_052_exp3_p060_413_bf0_obstruction_diagram.tex}
\end{center}

Consequently,
\begin{equation*}
 \tag{3}\label{sga3:iii:proposition:4-13:equation:3}
 b_{f_0}(c(X,Y,g\circ i'))(\bar x\otimes1)
 =\varphi\left(\sum_i\lambda_i\pi'(a'_i)\otimes1\right)
 =\sum_i\lambda_i f(a'_i),
\end{equation*}
the last equality following from \((\dagger)\). Comparison of the three
explicit formulas \SGARefLink{sga3:iii:proposition:4-13:equation:1}{(1)}, \SGARefLink{sga3:iii:proposition:4-13:equation:2}{(2)}, and \SGARefLink{sga3:iii:proposition:4-13:equation:3}{(3)} gives the required formula.

\medskip
\noindent\SGAInlineRefTarget{sga3:iii:corollary:4-14}\textbf{Corollary 4.14.}
\textit{Let \(Y,Y'\) be two flat subschemes of \(X\), both reducing to
\(Y_{\mathscr J}\), and suppose that \(Y_0\) is locally complete
intersection in \(X_0\). If \(f:T\to X\) is an \(S\)-morphism such that
\(f_{\mathscr J}\) factors through \(Y_{\mathscr J}\to X_{\mathscr J}\),
then}
\[
 c(X,Y,f)-c(X,Y',f)=b_{f_0}\bigl(d(Y,Y')\bigr).
\]

Indeed, applying the preceding formula to the diagram

\begin{center}
  \SGAThreeNativeInput[0.62\textwidth]
    {native_diagrams/exp3/diagram_053_exp3_p060_414_comparison_diagram.tex}
\end{center}

we find
\(c(X,Y,f)-c(X,Y',f)=b_{f_0}(c(X,Y,i'))\). Moreover, by \SGARefLink{sga3:iii:proposition:4-8:item:ii}{4.8(ii)},
\(c(X,Y,i')=d(Y,Y')\).

\medskip
\noindent\SGAInlineRefTarget{sga3:iii:proposition:4-15}\textbf{Proposition 4.15.}
\textit{Let \(X\) be an \(S\)-group smooth over \(S\), and let \(Y\)
be a sub-\(S\)-group flat and locally of finite presentation over \(S\).
Then \(Y\) is locally complete intersection in \(X\) (cf. \SGARefLink{sga3:iii:definition:4-6-2}{4.6.2}).}

\noindent\textit{Proof.}
\footnote{Editor's note: We have added to the statement the hypothesis
that \(Y\) be locally of finite presentation over \(S\), and supplied
the proof that follows, which is more direct than the one sketched in
the original. For completeness, we also spell out the latter. As in
the proof above, first reduce to the case \(S=\operatorname{Spec}(k)\),
where \(k\) is algebraically closed. By EGA IV\textsubscript{4},
16.9.10 and 19.3.2, it suffices to see that, for every \(y\in Y\), the
completion of the local ring \(\mathscr O_{Y,y}\) is a quotient of a
complete noetherian local ring by a regular sequence. By loc. cit.,
19.3.3, the set of \(y\in Y\) with this property is an open subset
\(U\subset Y\). Since \(Y\) is of finite type over \(k\), it suffices
to show that \(U\) contains every closed point. Because \(Y\) is a
\(k\)-group, a translation argument reduces this to the completion of
\(\mathscr O_{Y,e}\), that is, to the formal group \(\widehat Y\)
corresponding to \(Y\) (cf. Expos\'e VII\textsubscript{B}). Since \(X\)
is smooth, the affine algebra \(\mathscr A(\widehat X)\) is a power
series algebra \(k[[X_1,\ldots,X_n]]\). The Dieudonn\'e structure
theorem then shows that \(\mathscr A(\widehat Y)\) is isomorphic to a
quotient
\[
 k[[X_1,\ldots,X_{r+s}]]/
 (X_1^{p^{n_1}},\ldots,X_r^{p^{n_r}}),
\]
where \(p\) is the characteristic exponent of \(k\) and \(r+s\leq n\);
cf. Expos\'e VII\textsubscript{B}, Remark 5.5.2(b).}
We show that the immersion \(Y\to X\) is regular in the sense of EGA
IV\textsubscript{4}, 16.9.2. By EGA IV\textsubscript{4}, 19.5.1, this
implies that it is regular also in the sense of 4.6.2 (indeed, by loc.
cit., the two definitions are equivalent when \(S\) is locally
noetherian). Thus \lq\lq regular immersion\rq\rq\ will here mean regular
in the sense of EGA IV\textsubscript{4}, 16.9.2. Since \(X\) and \(Y\)
are flat and locally of finite presentation over \(S\), EGA
IV\textsubscript{4}, 19.2.4 reduces the assertion to proving that
\(Y_s\to X_s\) is a regular immersion for every \(s\in S\). By EGA
IV\textsubscript{4}, 19.1.5(ii), it is enough to check this on the
geometric fibers of \(S\), and hence when \(S\) is the spectrum of an
algebraically closed field \(k\).

By Expos\'e VI\textsubscript{A}, \SGARefLink{thm:sga3-VIA-3.2}{3.2}, the quotient \(X/Y\) exists and is
smooth, the morphism \(\pi:X\to X/Y\) is flat, and the following square
is cartesian:

\begin{center}
  \SGAThreeNativeInput[0.35\textwidth]
    {native_diagrams/exp3/diagram_054_exp3_p061_415_cartesian_square.tex}
\end{center}

Here \(\bar e\) is the image in \(X/Y\) of the unit point of \(X\).
By flat base change (cf. EGA IV\textsubscript{4}, 19.1.5(iii)), it
therefore suffices to show that \(i\) is a regular immersion. This is
immediate: the noetherian local ring \(\mathscr O_{X/Y,\bar e}\) is
smooth (hence regular), so its maximal ideal is generated by a regular sequence.

\medskip
\noindent\SGAInlineRefTarget{sga3:iii:paragraph:4-16}\textbf{4.16.}
\footnote{Editor's note: We have reorganized \SGARefLink{sga3:iii:paragraph:4-16}{4.16} by collecting here,
on the one hand, the hypotheses stated at the end of \SGARefLink{sga3:iii:proposition:4-15}{4.15} and, on the
other, the definition of the obstruction \(DY\).}
Let \(X\) be an \(S\)-group smooth over \(S\), and denote its group law
by \(\mu:X\times_S X\to X\). Let \(Y_{\mathscr J}\) be a
sub-\(S_{\mathscr J}\)-group of \(X_{\mathscr J}\), flat and locally of
finite presentation over \(S_{\mathscr J}\). By \SGARefLink{sga3:iii:proposition:4-15}{4.15},
\(Y_{\mathscr J}\) is locally complete intersection in \(X_{\mathscr J}\). Hence,
by \SGARefLink{sga3:iii:proposition:4-6-5}{4.6.5}, every flat
\(S\)-scheme \(Y\) lifting \(Y_{\mathscr J}\) is locally complete
intersection in \(X\). For such a \(Y\), \SGARefLink{sga3:iii:remark:4-8-0}{4.8.0} gives
\begin{equation*}
 \tag{4.16.1}\label{sga3:iii:paragraph:4-16:equation:4-16-1}
 \mathscr N_{Y/X}\otimes_{\mathscr O_Y}\mathscr O_{Y_0}
 =\mathscr N_{Y_0/X_0}
 =\mathscr N_{Y_{\mathscr J}/X_{\mathscr J}}
   \otimes_{\mathscr O_{Y_{\mathscr J}}}\mathscr O_{Y_0}.
\end{equation*}

On the other hand, let \(\varepsilon_0:S_0\to Y_0\) be the unit
section of \(Y_0\), and let \(\mathbf n_{Y_0/X_0}\) be the
quasi-coherent \(\mathscr O_{S_0}\)-module
\[
\mathbf n_{Y_0/X_0}=\varepsilon_0^*(\mathscr N_{Y_0/X_0}).
\]
Since \(Y_0\) and \(X_0\) are \(S_0\)-groups, it is immediate that
\(\mathscr N_{Y_0/X_0}\) is invariant under, say, left translations by
\(Y_0\). Hence
it is the inverse image of \(\mathbf n_{Y_0/X_0}\) under
\(Y_0\to S_0\); that is,
\begin{equation*}
 \tag{4.16.2}\label{sga3:iii:paragraph:4-16:equation:4-16-2}
 \mathscr N_{Y_0/X_0}
=\mathbf n_{Y_0/X_0}\otimes_{\mathscr O_{S_0}}\mathscr O_{Y_0}.
\end{equation*}

Taking (\SGARefLink{sga3:iii:paragraph:4-16:equation:4-16-1}{4.16.1}) and (\SGARefLink{sga3:iii:paragraph:4-16:equation:4-16-2}{4.16.2}) into account, we deduce first from \SGARefLink{sga3:iii:proposition:4-5}{4.5} that
the set of subschemes \(Y\) of \(X\), flat over \(S\) and lifting
\(Y_{\mathscr J}\), is either empty or principal homogeneous under
\begin{equation*}
 \tag{4.16.3}\label{sga3:iii:paragraph:4-16:equation:4-16-3}
 \operatorname{Hom}_{\mathscr O_{Y_0}}
\bigl(\mathbf n_{Y_0/X_0}\otimes_{\mathscr O_{S_0}}\mathscr O_{Y_0},
       \mathscr J\otimes_{\mathscr O_{S_0}}\mathscr O_{Y_0}\bigr).
\end{equation*}
On the other hand, \SGARefLink{sga3:iii:proposition:4-8:item:i}{4.8(i)} shows that, for every such \(Y\) and every
\(S\)-morphism \(f:T\to X\) such that
\(f_{\mathscr J}:T_{\mathscr J}\to X_{\mathscr J}\) factors through
\(Y_{\mathscr J}\), the obstruction \(c(X,Y,f)\) to \(f\) factoring
through \(Y\) belongs to
\[
 \operatorname{Hom}_{\mathscr O_{T_0}}
\bigl(\mathbf n_{Y_0/X_0}\otimes_{\mathscr O_{S_0}}\mathscr O_{T_0},
       \mathscr J\mathscr O_T\bigr).
\]
If \(T\) is also flat over \(S\), this last group is equal to
\[
 \operatorname{Hom}_{\mathscr O_{T_0}}
\bigl(\mathbf n_{Y_0/X_0}\otimes_{\mathscr O_{S_0}}\mathscr O_{T_0},
       \mathscr J\otimes_{\mathscr O_{S_0}}\mathscr O_{T_0}\bigr).
\]
This leads us to introduce the following group functor \(N_0\).

\medskip
\noindent\SGAInlineRefTarget{sga3:iii:definition:4-16-1}\textbf{Definition 4.16.1.}
\textit{Let \(N_0\) be the \(S_0\)-functor in commutative groups
defined, for every \(Z\in\operatorname{Ob}(\mathbf{Sch})_{/S_0}\), by}
\begin{equation*}
 \tag{*}\label{sga3:iii:definition:4-16-1:equation:star}
 \operatorname{Hom}_{S_0}(Z,N_0)
 =\operatorname{Hom}_{\mathscr O_Z}
\bigl(\mathbf n_{Y_0/X_0}\otimes_{\mathscr O_{S_0}}\mathscr O_Z,
       \mathscr J\otimes_{\mathscr O_{S_0}}\mathscr O_Z\bigr).
\end{equation*}
\textit{Then the set of subschemes \(Y\) of \(X\), flat over \(S\) and
lifting \(Y_{\mathscr J}\), is either empty or principal homogeneous
under}
\[
 \operatorname{Hom}_{S_0}(Y_0,N_0)=C^1(Y_0,N_0).
\]

For every such \(Y\), consider the diagram

\begin{center}
  \SGAThreeNativeInput[0.52\textwidth]
    {native_diagrams/exp3/diagram_055_exp3_p062_4161_DY_diagram.tex}
\end{center}

and let
\[
 DY=c\bigl(X,Y,\mu\circ(i,i)\bigr)
\]
be the obstruction to \(\mu\circ(i,i)\) factoring through \(Y\), that
is, to \(Y\) being stable under the group law of \(X\). By what
precedes, \(DY\) is an element of
\[
 N_0(Y_0\times_{S_0}Y_0)=C^2(Y_0,N_0).
\]
